\chapter{Using notebooks}
\label{\detokenize{tools-for-mathematics/01-using-notebooks/introduction/main:using-notebooks}}\label{\detokenize{tools-for-mathematics/01-using-notebooks/introduction/main:id1}}\label{\detokenize{tools-for-mathematics/01-using-notebooks/introduction/main::doc}}

\section{Introduction}
%\label{\detokenize{tools-for-mathematics/01-using-notebooks/introduction/main:introduction}}
At the advent of Calculus two mathematicians are credited with its
formalisation/invention:
\begin{itemize}
\item Isaac Newton
\item Gottfried Leibniz
\end{itemize}

One of the differences between the approaches taken by Newton and Leibniz is
their notation:

Newton denoted the derivative of a function \(f\) as:

\begin{equation*}
Df
\end{equation*}

and Leibniz denoted the derivative with the now more commonly used notation:
\begin{equation*}
\frac{df}{dx}
\end{equation*}

The mathematics itself is unchanged: what changes is the language/notation used
to communicate it.

Similarly when giving instructions through code to a computer there are a number
of notations, more commonly called languages available.

We will be using a language called \textbf{Python}.

Python was originally designed as a teaching language but it is now popular both
in academia and in industry.

In this chapter we will cover:
\begin{itemize}
\item 
Installing the specific distribution of Python on your computer.

\item 
Using something called a Jupyter notebook to write and run Python code.

\item Writing descriptive notes using \mintinline{markdown}{markdown} and \LaTeX\; (pronounced
    Lay-tech).

\end{itemize}

\section{Tutorial}
%\label{\detokenize{tools-for-mathematics/01-using-notebooks/tutorial/main:tutorial}}\label{\detokenize{tools-for-mathematics/01-using-notebooks/tutorial/main::doc}}

\subsection{Installation}
%\label{\detokenize{tools-for-mathematics/01-using-notebooks/tutorial/main:installation}}
\begin{enumerate}
\item Navigate to \texttt{https://www.anaconda.com/products/individual}.

\item Identify and download the version of Python 3 for your operating system
(Windows, MacOS, Linux).

\item Run the installer. I recommend using the default choices during the
installation process.

\end{enumerate}

If you have already used Python it is still recommended that you use the
Anaconda distribution. An explanation for this is available later.


\subsection{Starting a Jupyter notebook server}
%\label{\detokenize{tools-for-mathematics/01-using-notebooks/tutorial/main:starting-a-jupyter-notebook-server}}
We are going to use \textbf{Jupyter notebooks} for the first part of this book.
This interface to Python works inside your web browser but does not require an
internet connect.

Open a command line tool:
\begin{enumerate}
\item 
On \textbf{Windows} search for \texttt{Anaconda Prompt} (it should be available to you
after installing Anaconda). See Figure~\ref{fig:starting-the-notebook-server-windows}.

\item On \textbf{MacOS} search for \texttt{terminal}. See
See Figure~\ref{fig:starting-the-notebook-server}.


\end{enumerate}

In there type:

\begin{cliin}
$ jupyter notebook
\end{cliin}

and then press \texttt{Enter} on your keyboard.

Throughout this book, when there are commands to be typed in a command line
tool I will prefix them with a \texttt{\$}. Do not type the \texttt{\$}.

\begin{figure}[htbp]
\centering

    \includegraphics[width=0.750\linewidth]{assets/starting_the_notebook_server/main.png}
    \caption{Starting the notebook server on
    MacOS}\label{fig:starting-the-notebook-server}
\end{figure}

\begin{figure}[htbp]
\centering

    \includegraphics[width=0.750\linewidth]{assets/starting_the_notebook_server_windows/main.png}
    \caption{Starting the notebook server on
    Windows}\label{fig:starting-the-notebook-server-windows}
\end{figure}

This will open a new page in your browser. The url bar at the top should have
something that looks like: \texttt{http://localhost:8888/tree}.


\begin{figure}[htbp]
\centering
    \includegraphics[width=0.750\linewidth]{assets/the_jupyter_interface/main.png}
    \caption{The Jupyter notebook interface}
\end{figure}

This is the general interface to the Jupyter server. It shows the general file
structure on your computer.


\subsection{Creating a new notebook}
%\label{\detokenize{tools-for-mathematics/01-using-notebooks/tutorial/main:creating-a-new-notebook}}
In the top right, click on the \texttt{new} button and click on \texttt{Python 3}.

\begin{figure}[htbp]
\centering
    \includegraphics[width=0.750\linewidth]{assets/creating_a_new_notebook/main.png}
    \caption{Creating a new notebook}
\end{figure}

Let us change the name of the notebook by clicking on “Untitled” and changing
the name. We will call it “introduction”.

\begin{figure}[htbp]
\centering
    \includegraphics[width=0.750\linewidth]{assets/changing_notebook_name/main.png}
    \caption{Changing the notebook name}
\end{figure}

Once this is done let us close the notebook by closing the corresponding tab of your web browser.

\subsection{Organising our files}
%\label{\detokenize{tools-for-mathematics/01-using-notebooks/tutorial/main:organising-our-files}}

Open your file browser:
\begin{enumerate}
    \item File Explorer on \textbf{Windows} (see
        Figure~\ref{fig:new_directory_structure_windows})

\item Finder on \textbf{MacOS} (see
        Figure~\ref{fig:new_directory_structure})
\end{enumerate}

Navigate to where your notebook is (this might not be immediately evident): you
should see a \texttt{introduction.ipynb} file.

Let us rearrange things.

Find a location on your computer where you want to keep the files for this
book, using your file browser:

\begin{enumerate}
\item Create a new directory called \texttt{cfm} (short for “Computing for Mathematics”);

\item Inside that directory create a new directory called \texttt{nbs} (short for
``Notebooks'');

\item Move the \texttt{introduction.ipynb} file to this \texttt{nbs} directory.

\end{enumerate}

\begin{figure}[htbp]
\centering
    \includegraphics[width=0.750\linewidth]{assets/new_directory_structure/main.png}
    \caption{Creating a new directory on MacOS}\label{fig:new_directory_structure}
\end{figure}


\begin{figure}[htbp]
\centering
    \includegraphics[width=0.750\linewidth]{assets/new_directory_structure_windows/main.png}
    \caption{Creating a new directory on
    Windows}\label{fig:new_directory_structure_windows}
\end{figure}


\subsection{Writing some basic Python code}
%\label{\detokenize{tools-for-mathematics/01-using-notebooks/tutorial/main:writing-some-basic-python-code}}
Go back to our Jupyter notebook server (in your browser).

Use the interface to navigate to the \texttt{cfm} directory and inside that the \texttt{nbs}
directory and open the \texttt{introduction.ipynb} notebook.

\begin{figure}[htbp]
\centering
    \includegraphics[width=0.750\linewidth]{assets/opening_notebook/main.png}
    \caption{Opening a notebook}
\end{figure}

In the first available “cell” write the following calculation:

\begin{pyin}
    2 + 2
\end{pyin}


When you have done that click on the \texttt{Run} button (you can also use
\texttt{Shift + Enter} as a keyboard shortcut).

\begin{pyin}
    2 + 2
\end{pyin}

\[2\]

We see two different things there:

\begin{enumerate}
\item 

The input: \texttt{In {[}1{]}} which is the instruction to Python to use the
mathematical technique of addition to compute 2 + 2.

\item 

The output: \texttt{Out {[}1{]}} showing the output that Python has returned as a result
of the instruction.

\end{enumerate}


\subsection{Writing markdown}
%\label{\detokenize{tools-for-mathematics/01-using-notebooks/tutorial/main:writing-markdown}}

One of the reasons for using Jupyter notebooks is that it allows us to include
both code and communication using something called \texttt{markdown}.


Create a new cell and change the cell type to \texttt{Markdown}. Now write the
following in there:

\begin{minted}{markdown}
As well as using Python in Jupyter notebooks we can also write using Markdown.
This allows us to use basic \LaTeX\; as a way to display mathematics.
For example:

1. $\frac{2}{3}$
2. $\sum_{i=0}^n i$
\end{minted}


When you run that it should look like Figure~\ref{fig:rendering_markdown}.

\begin{figure}[htbp]
\centering
    \includegraphics[width=0.750\linewidth]{assets/rendering_markdown/main.png}
    \caption{Rendering markdown}\label{fig:rendering_markdown}
\end{figure}

\subsection{Saving your notebook to a different format}
%\label{\detokenize{tools-for-mathematics/01-using-notebooks/tutorial/main:saving-your-notebook-to-a-different-format}}

Click on \texttt{File} and \texttt{Download As} this brings up a number of formats that
Jupyter notebooks can be exported to. Some of these might need other tools
installed on your computer but a portable option is \texttt{HTML}.


Click on \texttt{HTML (.html)}.


Now use your file browser and open the downloaded file. This will open in your
browser a static version of the file you have been working on. This is a helpful
way to share your work with someone who might not have Jupyter (or even Python).

In this tutorial we have
\begin{itemize}
\item 

Installed the Anaconda distribution of Python.

\item 

Started a notebook server.

\item 

Created a new notebook.

\item 

Run some Python code.

\item 

Written some markdown.

\item 

Saved our notebook to a different format.

\end{itemize}




\section{How}
%\label{\detokenize{tools-for-mathematics/01-using-notebooks/how/main:how}}\label{\detokenize{tools-for-mathematics/01-using-notebooks/how/main::doc}}

\subsection{Install Anaconda}
%\label{\detokenize{tools-for-mathematics/01-using-notebooks/how/main:install-anaconda}}

\begin{enumerate}

\item 

Navigate to \texttt{https://www.anaconda.com/products/individual}

\item 

Download the distribution of anaconda for your Operating System

\item 

Run the installer

\end{enumerate}


\subsection{Start a Jupyter notebook server}
%\label{\detokenize{tools-for-mathematics/01-using-notebooks/how/main:start-a-jupyter-notebook-server}}


\begin{enumerate}

\item 

    Open a command line tool (\texttt{Anaconda prompt}) on Windows,
        \texttt{terminal} on OS X);

\item 

    Type \texttt{jupyter notebook} and press \texttt{Enter}

\end{enumerate}


\subsection{Create a new notebook}
%\label{\detokenize{tools-for-mathematics/01-using-notebooks/how/main:create-a-new-notebook}}
\begin{enumerate}

\item 

Navigate to the location you want using the Jupyter interface;

\item 

Click on the \texttt{new} button in the top right;

\item 

Rename the notebook (change \texttt{Untitled} in the top left to a name of your choice).

\end{enumerate}


\subsection{Find/open a notebook}
%\label{\detokenize{tools-for-mathematics/01-using-notebooks/how/main:find-open-a-notebook}}

Using a file browser you can navigate the directories and files on your
computer. Jupyter notebooks appear as generic files with the \texttt{.ipynb}
extension.


You cannot double click on these to open them, you need to navigate to them through the Jupyter interface.


\subsection{Run Python code}
%\label{\detokenize{tools-for-mathematics/01-using-notebooks/how/main:run-python-code}}

In a Jupyter notebook cell write an instruction, for example:

\begin{pyin}
3 / 5
\end{pyin}


and click on the \texttt{Run} button or use \texttt{Shift + Enter} as a keyboard shortcut.


\subsection{Carry out basic arithmetic operations}
%\label{\detokenize{tools-for-mathematics/01-using-notebooks/how/main:carry-out-basic-arithmetic-operations}}\label{\detokenize{tools-for-mathematics/01-using-notebooks/how/main:id1}}

The python code for the following arithmetic operations are:

\begin{enumerate}

\item 

Addition, \(2 + 2\): \texttt{2 + 2};

\item 

Subtraction, \(3 - 1\): \texttt{3 - 1};

\item 

Multiplication, \(3 \times 5\): \texttt{3 * 5};

\item 

Division, \(20 / 5\): \texttt{20 / 5};

\item 

Exponentiation, \(2 ^ 4\): \texttt{2 ** 4};

\item 

Integer remainder, \(5 \mod 2\): \texttt{5 \% 2};

\item 

Combining operations, \(\frac{2 ^ 3 + 1}{4}\): \texttt{(2 ** 3 + 1) / 4};

\end{enumerate}


\textbf{Note} that instructions to a computer (through the code we write) need
to be specific. For example the \texttt{\textasciicircum{}} symbol in Python
does not mean exponentiation. If you were to type \texttt{2 \textasciicircum{} 4} you would get an error.

In later chapters we will see what the specific instructions are to carry out more complex operations.

\subsection{Write markdown}
%\label{\detokenize{tools-for-mathematics/01-using-notebooks/how/main:write-markdown}}

To write markdown click on a cell and change the type to \texttt{Markdown}, you can do
this by click on \texttt{Cell}, \texttt{Cell Type} or by using the scroll wheel in the menu
bar.


Markdown is a lightweight ``mark up'' language that allows you to write and
include various types of formatting which include:

\begin{enumerate}

\item 

Headings;

\item 

Bold and italics;

\item 

Ordered and unordered lists;

\item 

Code (which will only be displayed but not run);

\item 

Hyperlinks

\end{enumerate}


The syntax is relatively straightforward and the following is a good guide:


\texttt{https://www.markdownguide.org/cheat-sheet/}


\subsection{Write basic LaTeX}
%\label{\detokenize{tools-for-mathematics/01-using-notebooks/how/main:write-basic-latex}}

Jupyter notebooks allow for markdown cells to not only include markdown but also
include mathematics using another ``mark up''' language called \LaTeX.


Here is a brief overview of the syntax for arithmetic operations:
\begin{itemize}
\item 

\texttt{\$a+b\$} gives: \(a + b\):

\item 

\texttt{\$a-b\$} gives: \(a-b\)

\item 

\texttt{\$-a\$} gives: \(-a\)

\item 

\texttt{\$ab\$} gives \(ab\)

\item 

\texttt{\$a\textbackslash{}cdot b\$} gives \(a\cdot b\)

\item 

\texttt{\$a\textbackslash{}times b\$} gives \(a\times b\)

\item 

\texttt{\$a/b\$} gives \(a/b\)

\item 

\texttt{\$\textbackslash{}frac\{a\}\{b\}\$} gives \(\frac{a}{b}\)

\item 

\texttt{\$a \textasciicircum{} b\$} gives \(a ^ b\)

\end{itemize}


The \texttt{\$<expression>\$} delimiters create what is called an ``inline'''
mathematics environment.
You can change these to \texttt{\$\$<expression>\$\$} to give ``displayed
mathematics''.


We can write a matrix:

\begin{minted}{markdown}
    \begin{pmatrix}
        a&b\\
        c&d\\
        e&f\\
    \end{pmatrix}
\end{minted}


gives:

\begin{equation*}
\begin{split}
    \begin{pmatrix}
        a&b\\
        c&d\\
        e&f\\
    \end{pmatrix}
\end{split}
\end{equation*}

We can write integrals:

\begin{minted}{markdown}
$$
    \int_{0}^{\infty}x dx
$$
\end{minted}


gives:

\begin{equation*}
\begin{split}
    \int_{0}^{\infty}x dx
\end{split}
\end{equation*}

We can write summations:

\begin{minted}{python}
$$
    \sum_{0}^{n}i
$$
\end{minted}


gives:

\begin{equation*}
\begin{split}
    \sum_{0}^{n}i
\end{split}
\end{equation*}

\subsection{Save the output in a different format}
%\label{\detokenize{tools-for-mathematics/01-using-notebooks/how/main:save-the-output-in-a-different-format}}

Click on \texttt{File} then \texttt{Download as} and choose the format you want
to use. \texttt{HTML}
is a portable option that can be viewed on most devices, note however that you
cannot run the cells: what you are downloading a is static version of your
notebook.

\section{Exercises}
%\label{\detokenize{tools-for-mathematics/01-using-notebooks/exercises/main:exercises}}\label{\detokenize{tools-for-mathematics/01-using-notebooks/exercises/main::doc}}

\textbf{If you are not sure how to do something, have a look at the “How To” section.}

\begin{enumerate}

\item 

Create a new notebook rename it “exercises”. Navigate to it using your file browser to make sure you can find it.

\item 

Write and run some Python code to carry out the following calculations
\begin{enumerate}

\item 

\(3 + 8\)

\item 

\(3 / 7\)

\item 

\(456 / 21\)

\item 

\(\frac{4 ^ 3 + 2}{2\times 5} - 5 ^ {\frac{1}{2}}\)

\end{enumerate}

\item 

Write a markdown cell with the following and view the rendered version:

\begin{minted}{markdown}
# Euler's equation

$$
e ^ {i\pi} = -1
$$
\end{minted}

\item 

Render the following expressions by writing markdown
\begin{enumerate}

\item 

\(\frac{4 ^ 3 + 2}{2\times 5}\)

\item 

\(- 5 ^ {\frac{1}{2}}\)

\item 

\(\frac{df}{dx}\)

\item 

\(\int_{5}^{12}x^2dx\)

\item 

\(\begin{pmatrix}4 & 12 & 3\\2 & x & i\\\end{pmatrix}\)

\end{enumerate}

\item 

Save your notebook to \texttt{HTML} and open and view it.

\item 

Download the notebook available at
\texttt{10.5281/zenodo.7118738} and
check that you are able to open it.

\end{enumerate}




\section{Solutions}
%\label{\detokenize{tools-for-mathematics/01-using-notebooks/solutions/main:solutions}}\label{\detokenize{tools-for-mathematics/01-using-notebooks/solutions/main::doc}}

\label{\detokenize{tools-for-mathematics/01-using-notebooks/solutions/main:question-1}}

\begin{enumerate}
    \item  Figure~\ref{fig:create_a_new_notebook} shows the steps.

        \begin{figure}[!htbp]
                \centering
                \begin{subfigure}[b]{0.4\textwidth}
                    \includegraphics[width=\textwidth]{assets/create_a_new_notebook/main00050.png}
                    \caption{Creating a blank notebook}
                    \label{fig:create_a_blank_notebook}
                \end{subfigure}
                \hfill
                \begin{subfigure}[b]{0.4\textwidth}
                    \includegraphics[width=\textwidth]{assets/create_a_new_notebook/main00069.png}
                    \caption{Renaming the Untitled notebook}
                    \label{fig:renaming_the_untitiled_notebook}
                \end{subfigure}
                \hfill
                \begin{subfigure}[b]{0.4\textwidth}
                    \includegraphics[width=\textwidth]{assets/create_a_new_notebook/main00138.png}
                    \caption{Saving the new name}
                    \label{fig:saving_the_new_name}
                \end{subfigure}
                \hfill
                    \begin{subfigure}[b]{0.4\textwidth}
                    \includegraphics[width=\textwidth]{assets/create_a_new_notebook/main00239.png}
                    \caption{Viewing the new notebook}
                    \label{fig:viewing_the_new_notebook}
                \end{subfigure}
                        
                \caption{Creating a new notebook}
                \label{fig:create_a_new_notebook}
        \end{figure}

    \item 

        \begin{enumerate}

            \item \(3 + 8\)
            \begin{pyin}
3 + 8
                \end{pyin}
                \[11\]

            \item \(3 / 7\)
            \begin{pyin}
3 / 7
                \end{pyin}
                \[0.42857142857142855\]

            \item \(456 / 21\)
            \begin{pyin}
456 / 21
                \end{pyin}
                \[21.714285714285715\]

            \item \(\frac{4 ^ 3 + 2}{2\times 5} - 5 ^ {\frac{1}{2}}\)

            \begin{pyin}
(4 ** 3 + 2) / (2 * 5) - 5 ** (1 / 2)
                \end{pyin}
                \[4.36393202250021\]
        \end{enumerate}

    \item 

    \begin{minted}{markdown}
# Euler's equation

$$
e ^ {i\pi} = -1
$$
    \end{minted}

    Figure~\ref{fig:writing_markdown} shows the markdown cell and the rendered
    output.

        \begin{figure}[!htbp]
                \centering
                \begin{subfigure}[b]{0.4\textwidth}
                    \includegraphics[width=\textwidth]{assets/writing_markdown/main00295.png}
                    \caption{Writing in a markdown cell}
                    \label{fig:writing_in_a_markdown_cell}
                \end{subfigure}
                \hfill
                \begin{subfigure}[b]{0.4\textwidth}
                    \includegraphics[width=\textwidth]{assets/writing_markdown/main00342.png}
                    \caption{The rendered output}
                    \label{fig:the_rendered_output}
                \end{subfigure}
                \caption{Writing markdown}
                \label{fig:writing_markdown}
        \end{figure}

    \item 

        \begin{enumerate}
            \item 
                \begin{minted}{markdown}
$\frac{4 ^ 3 + 2}{2\times 5}$
                \end{minted}
                    \[\frac{4 ^ 3 + 2}{2\times 5}\]
            \item 
                \begin{minted}{markdown}
$- 5 ^ {\frac{1}{2}}$
                \end{minted}
                    \[- 5 ^ {\frac{1}{2}}\]
            \item 
                \begin{minted}{markdown}
$\frac{df}{dx}$
                \end{minted}
                    \[\frac{df}{dx}\]
            \item 
                \begin{minted}{markdown}
$\int_{5}^{12}x^2dx$
                \end{minted}
                    \[\int_{5}^{12}x^2dx\]
            \item 
                \begin{minted}{markdown}
$\begin{pmatrix}4 & 12 & 3\\2 & x & i\\\end{pmatrix}$
                \end{minted}
                    \[\begin{pmatrix}4 & 12 & 3\\2 & x & i\\\end{pmatrix}\]
        \end{enumerate}

    \item Figure~\ref{fig:saving_and_opening_html} shows the steps.


        \begin{figure}[!htbp]
                \centering
                \begin{subfigure}[b]{0.4\textwidth}
                    \includegraphics[width=\textwidth]{assets/saving_and_opening_html/main00055.png}
                    \caption{Saving the notebook as html}
                    \label{fig:saving_the_notebook_as_html}
                \end{subfigure}
                \hfill
                \begin{subfigure}[b]{0.4\textwidth}
                    \includegraphics[width=\textwidth]{assets/saving_and_opening_html/main00210.png}
                    \caption{Viewing the html file}
                    \label{fig:viewing_the_html_file}
                \end{subfigure}
                \caption{Saving and opening the html file}
                \label{fig:saving_and_opening_html}
        \end{figure}

    \item Figure~\ref{fig:downloading_a_notebook} shows the steps.


        \begin{figure}[!htbp]
                \centering
                \begin{subfigure}[b]{0.3\textwidth}
                    \includegraphics[width=\textwidth]{assets/downloading_a_notebook/main00150.png}
                    \caption{Downloading the notebook file to the Desktop
                    directory}
                    \label{fig:downloading_the_notebook_file_to_the_desktop}
                \end{subfigure}
                \hfill
                \begin{subfigure}[b]{0.3\textwidth}
                    \includegraphics[width=\textwidth]{assets/downloading_a_notebook/main00330.png}
                    \caption{Navigating to the notebook file in the Jupyter
                    interface}
                    \label{fig:navigating_to_the_notebook_file_in_the_jupyter_interface}
                \end{subfigure}
                \hfill
                \begin{subfigure}[b]{0.3\textwidth}
                    \includegraphics[width=\textwidth]{assets/downloading_a_notebook/main00498.png}
                    \caption{Viewing the notebook}
                    \label{fig:viewing_the_notebook}
                \end{subfigure}
                \caption{Saving and opening the html file}
                \label{fig:downloading_a_notebook}
        \end{figure}

\end{enumerate}

\section{Further information}
%\label{\detokenize{tools-for-mathematics/01-using-notebooks/why/main:further-information}}\label{\detokenize{tools-for-mathematics/01-using-notebooks/why/main::doc}}

\subsection{Why use anaconda?}
%$\label{\detokenize{tools-for-mathematics/01-using-notebooks/why/main:why-use-anaconda}}

Python is a free an open source piece of software. One of the main reasons for
its popularity is that there are a number of separate tools that work well with
it, these are called libraries. Sometimes installing these libraries can require
an understanding of some potential pitfalls. In scientific circles the Anaconda
distribution was developed to give a single download of not only Python but a
lot of commonly used libraries.


If you want to read more about this here are some web resources:
\begin{itemize}
\item 

\texttt{https://medium.com/@RustyComer/why-use-anaconda-524bb6765e4d}

\item 

\texttt{https://www.reddit.com/r/Python/comments/betkoj/why\_use\_anaconda/}

\end{itemize}


\subsection{Why use Jupyter?}
%\label{\detokenize{tools-for-mathematics/01-using-notebooks/why/main:why-use-jupyter}}

There are are variety of ways to write and run Python:
\begin{enumerate}

\item 

Using an interactive notebook environment like Jupyter;

\item 

Using an integrated development environment and/or editor.

\end{enumerate}


We will in fact use the second approach in the second part of this book.


One strength of Jupyter are that it allows you to include communication (writing
through markdown) with your code. This allows you to use code and describe what
you are using it for.


Another advantage is that it allows you to immediately have your output next to your input.


There are some limitations to Jupyter as an editor which is why we will explore
using a powerful editor in the second part of the course.


In general:
\begin{enumerate}

\item 

Jupyter is a fantastic way to interactively use and communicate code;

\item 

Integrated development environments and/or editors are the correct tool to
write software.

\end{enumerate}


In this course you will learn to use either approach in the appropriate manner
for the right task. For the first part we will mainly be using code
interactively and so we will use Jupyter notebooks.


If you are interested here are some further resources:
\begin{itemize}
\item 

Why Jupyter Notebooks Are So Popular Among Data Scientists: 

        \texttt{https://analyticsindiamag.com/why-jupyter-notebooks-are-so-popular-among-data-scientists/}

\item 

10 reasons why data scientists love Jupyter notebooks: 

        \texttt{https://hub.packtpub.com/10-reasons-data-scientists-love-jupyter-notebooks/}

\item 

I don’t like notebooks - Joel Grus (Allen Institute for Artificial
        Intelligence): \texttt{https://www.youtube.com/watch?v=7jiPeIFXb6U}

\end{itemize}


Some further information on using Jupyter:
\begin{itemize}
\item 

\texttt{https://jupyter-notebook-beginner-guide.readthedocs.io/}

\item 

\texttt{https://www.analyticsvidhya.com/blog/2018/05/starters-guide-jupyter-notebook/}

\end{itemize}


\subsection{Why can I not double click on a Jupyter notebook file?}
%\label{\detokenize{tools-for-mathematics/01-using-notebooks/why/main:why-can-i-not-double-click-on-a-jupyter-notebook-file}}

When you double click on a file and your computer opens it in an application
that is because a default is set for the particular file extension. For example
double clicking on \texttt{main.docx} will automatically open up the document using a
word processor (like Microsoft word). This is because the file has the extension
\texttt{.docx} and your operating system has set that anything with that extension will
be opened in that particular application. You could also open the
application and navigate to the file and open it directly.


With Jupyter notebooks no default is set by the operating system as the
application that opens it is in fact a local web server in your browser. As such
you do not have a choice and need to open it in the Jupyter interface.


\subsection{Where can I find keyboard shortcuts for using Jupyter}
%\label{\detokenize{tools-for-mathematics/01-using-notebooks/why/main:where-can-i-find-keyboard-shortcuts-for-using-jupyter}}

In a notebook if you go to the menu bar and click on \texttt{Help} followed by
\texttt{Keyboard Shortcuts} you will find a number of helpful keyboard shortcuts.


For example, when on a cell pressing \texttt{Esc} followed by \texttt{m} will turn the cell in
to a markdown cell.


\subsection{What is markdown?}
%\label{\detokenize{tools-for-mathematics/01-using-notebooks/why/main:what-is-markdown}}

As described here \texttt{https://www.markdownguide.org/getting-started/}:

\begin{quote}
Markdown is a lightweight markup language that you can use to add formatting
elements to plain text text documents. Created by John Gruber in 2004,
Markdown is now one of the world’s most popular markup languages
\end{quote}


\subsection{What is LaTeX?}
%\label{\detokenize{tools-for-mathematics/01-using-notebooks/why/main:what-is-latex}}

As described here \texttt{https://www.latex-project.org/about/}:

\begin{quote}
LaTeX;, which is pronounced ``Lah-tech'' or ``Lay-tech'' (to rhyme with «blech» or
``Bertolt Brecht'''), is a document preparation system for high-quality
typesetting. It is most often used for medium-to-large technical or scientific
documents but it can be used for almost any form of publishing.
\end{quote}

\begin{quote}
LaTeX\; is not a word processor! Instead, LaTeX\; encourages authors not to worry
too much about the appearance of their documents but to concentrate on getting
the right content.
\end{quote}


We will learn more about \LaTeX\; in the later part of this course but for
now we only need to know that it we can use \LaTeX\; to write an instruction
for Jupyter to display mathematics.


\subsection{Can I use \texttt{\textbackslash(} and \texttt{\textbackslash)} instead of \texttt{\$} for \LaTeX?}
%\label{\detokenize{tools-for-mathematics/01-using-notebooks/why/main:can-i-use-and-instead-of-for-latex}}\label{\detokenize{tools-for-mathematics/01-using-notebooks/why/main:why-do-we-use-brackets-instead-of-dollars-for-latex}}

You will see in some places that \texttt{\textbackslash(},
\texttt{\textbackslash)} or \texttt{\textbackslash[},
\texttt{\textbackslash]} can be used as delimiters for
\LaTeX\; when used outside of Jupyter notebooks. This is in fact recommended for a
number of reasons one of which are given at
\texttt{https://vknight.org/tex/\#12-inline-mathematics}:

\begin{quote}
    Note that using \texttt{\textbackslash(} and \texttt{\textbackslash)} is preferred over \$. One of the reasons is that it
is easier for humans (and machines) to find the start and end of some
mathematics.
\end{quote}


If you want to use \texttt{\textbackslash(}, \texttt{\textbackslash)} or
\texttt{\textbackslash[}, \texttt{\textbackslash]} as
mathematics delimiters within Jupyter notebooks you need to escape the
\texttt{\textbackslash} and use: \texttt{\textbackslash\textbackslash(},
\texttt{\textbackslash\textbackslash)} or
\texttt{\textbackslash\textbackslash[},
\texttt{\textbackslash\textbackslash]} instead.


\subsection{Other resources for \LaTeX\; syntax.}
\label{\detokenize{tools-for-mathematics/01-using-notebooks/why/main:other-resources-for-latex-syntax}}

This cheat sheet: \texttt{http://tug.ctan.org/info/undergradmath/undergradmath.pdf}


This web page is my recommended set of resources for learning LaTeX:
\texttt{https://vknight.org/tex/}.


\subsection{What is a markup language?}
\label{\detokenize{tools-for-mathematics/01-using-notebooks/why/main:what-is-a-markup-language}}

\LaTeX\; and markdown are both examples of what is called a \textbf{markup language}.
Another common example of a markup language is html (the way web pages are
written).


A markup language is a system that allows us to write content alongside
annotations to specify how the content is to appear.


This description of markdown from
\texttt{https://www.markdownguide.org/getting-started/} is not specific to markdown but
to any markup language:
\begin{quote}


Using Markdown is different than using a WYSIWYG editor. In an application
like Microsoft Word, you click buttons to format words and phrases, and the
changes are visible immediately. Markdown isn’t like that. When you create a
Markdown-formatted file, you add Markdown syntax to the text to indicate which
words and phrases should look different.
\end{quote}


In general whilst it might take a little while to learn all the intricacies of a
markup language it allows for more portability and precision.


Markup languages differ in complexity:
\begin{itemize}
\item 

\LaTeX\; is incredibly sophisticated and has a huge range of capabilities.

\item 

Markdown is designed to be basic with a few specific annotations to remember.

\end{itemize}


\subsection{I only have a chromebook: how can I use notebooks?}
%\label{\detokenize{tools-for-mathematics/01-using-notebooks/why/main:i-only-have-a-chromebook-how-can-i-use-notebooks}}

There are a number of cloud based services that give access to scientific Python
environments. This 2018 article gives a brief review of 5 of them:

\texttt{https://medium.com/@siddesh.001/top-5-online-free-notebook-ipynb-and-other-cloud-services-dbf9580d99e3}

I recommend using cocalc: \texttt{https://cocalc.com}. The free tier does have some
limitations but it should be sufficient to be able to work through this book.


\subsection{I only have a tablet: how can I use notebooks?}
\label{\detokenize{tools-for-mathematics/01-using-notebooks/why/main:i-only-have-a-tablet-how-can-i-use-notebooks}}

There are two iOS apps that I am aware of for notebooks:
\begin{itemize}
\item 

\texttt{Carnets}

\item 

\texttt{Juno}

\end{itemize}
